\chapter{Atomic Energies and the Periodic Table}\label{ch:atoms}
\index{atomic energies}\index{shell structure}\index{sub-shell, (4+4)}\index{octet}\index{pair}%
\index{capacity}\index{parity rule}\index{Linnett double quartet}\index{Neumann boundary condition}%
\index{periodic table}\index{period doubling}

% Material to be added later from prel book (realquantum2.tex):
%   - chapters/sphericalsym.tex (computations in spherical symmetry — original
%     1D radial-equation treatment of atoms in the 2017 draft)

\begin{quote}
\itshape Would it not be pre-established harmony of a peculiar sort if the classical-epoch researchers, those who, as we hear today, had no idea of what measuring truly is, had unwittingly gone on to give us as legacy a guidance scheme revealing just what is fundamentally measurable for instance about a hydrogen atom?\\
\normalfont\hfill --- E.~Schr\"odinger, \emph{Die gegenw\"artige Situation in der Quantenmechanik}, 1935.
\end{quote}
\bigskip

Part~III opens with the Level-1 atomic model and its validation against observed total atomic
energies. If Level 1 does not reproduce these to acceptable accuracy, the rest of the hierarchy
has nothing to stand on. If it does, every higher-level architecture is anchored.

This chapter reports two things, and they are of different standing. The first is a
\emph{calculation}: neutral atoms computed all-electron from hydrogen to krypton, agreeing with
reference total energies to a few per cent, with the nuclear charge $Z$ as the only supplied
number. The second is a \emph{construction}: two filled packings appear and no third --- the
antipodal pair $2$ and the octet $8=4{+}4$ --- and the row lengths of the periodic table are sums
of these two alone. The construction is the more interesting of the two, and also the one with the
sharper limits: capacities are derived at every $Z$, the filling only as far as carbon and
supplied beyond.

\section{The Level-1 atomic model}\label{sec:level1-model}

The Level-1 model of Section~\ref{sec:level1} is: a point kernel of charge $+Z$ at the origin
surrounded by $N$ one-electron charge densities $\psi_i^2$ on non-overlapping domains $\Omega_i$,
continuous across free interfaces, interacting by Coulomb potentials without self-repulsion, with
the ground state reached by parabolic relaxation.

The only inputs are:
\begin{itemize}
\item the nuclear charge $Z$;
\item the electron number $N$, equal to $Z$ for a neutral atom;
\item a computational mesh.
\end{itemize}
There is no orbital basis, no exchange--correlation functional, no antisymmetry, no two-electron
integrals, and no choice of angular momentum to populate.

It is worth being exact about what ``no fitted parameters'' does and does not cover here, because
an earlier form of this claim was too broad. Every atom in this chapter is solved
\textbf{all-electron with the bare kernel $-Z/r$}: no core radius, no pseudo-kernel, no
two-scale core/valence decomposition, no fitted length, nothing adjusted per element. Those
devices are real and are used later in the book --- the pseudo-kernel reduction of
Chapter~\ref{ch:hydrides} is built on a calibrated $(Z_{\rm kernel}, r_c)$ --- but they are
reductions \emph{downstream} of this chapter, and the atomic anchor itself does not rest on them.

\section{Two numbers, not four: the pair and the octet}\label{sec:twoeight}

The shell structure is not imposed. What the packing admits are two filled arrangements, and the
reason there is no third is a symmetry argument rather than an energy comparison.

\paragraph{The primitives are $2$ and $4$.}
The antipodal pair and the tetrahedron. Eight is \emph{not} a third primitive: it is $4{+}4$, two
dual tetrahedra. In the atom the tetrahedron never closes a shell on its own.

\paragraph{Why the sequence stops at eight.}
A closed shell must be spherically featureless --- it must present no direction to the charge
arriving next. A lone tetrahedron has $T_d$ symmetry and \emph{no inversion centre}, so it is
directional and cannot close; it admits its own inversion, a genuinely different set of four, and
the union of the two is the cube, $O_h$, which \emph{is} centrosymmetric. The sequence stops
there, the cube being already its own inversion. Inversion is an involution, so there is exactly
one partner and never two. The pair is centrosymmetric immediately, so it closes at $2$ --- and
this is also why $2$ is the one capacity that is not internally doubled.

\paragraph{Row capacities, and a parity result.}
Row capacities are sums of the pair and the octet with nothing else entering:
\[
2,\qquad 8,\qquad 18 = 2{+}8{+}8,\qquad 32 = 8{+}8{+}8{+}8 .
\]
That these two units suffice is not an accident of the first few rows. Writing $c_n$ for the
capacity of the $n$th shell, $c_n \equiv 2 \pmod 8$ for odd $n$ and $c_n \equiv 0 \pmod 8$ for even
$n$: the remainder is only ever $0$ or $2$, and \emph{never} $4$ or $6$. That is the content of the
observation, since a row of length $8k{+}4$ or $8k{+}6$ could not be composed from
$\{2,8\}$ at all. So every capacity is a whole number of octets carrying at most one pair, the
pair appearing exactly when $n$ is odd. In closed form, entirely in the model's own units,
\[
\mathrm{cap}(n) \;=\; 8\,\Big\lfloor \tfrac n2 \Big\rfloor \Big\lceil \tfrac n2 \Big\rceil
\;+\; 2\,(n \bmod 2),
\]
the octet count being the product of the two halves of $n$, and a pair present if and only if $n$
is odd. This equals $2n^2$ for every $n$ --- but it never mentions $2n^2$, and that is the point:
the account does not borrow the hydrogenic degeneracy, it \emph{coincides} with it.

\paragraph{What follows, and what does not.}
Two claims must be kept apart, and conflating them was an error corrected during review.
\begin{enumerate}
\item[(A)] Each row length decomposes into pairs and octets, uniquely, with a pair present iff $n$
is odd. \textbf{This holds.}
\item[(B)] The row lengths \emph{follow from} the model. \textbf{This is not established.} The
configuration generator takes the row lengths as input; balanced filling then distributes charges
within capacities it has been given.
\end{enumerate}

\paragraph{Period doubling is open.}
Each capacity serves two periods --- $8,8$ then $18,18$ then $32,32$ --- and this is \emph{not}
explained by the symmetry argument above, which accounts for the internal doubling $4{+}4=8$ but
says nothing about why a capacity repeats. The cube is already centrosymmetric, so there is no
second involution to appeal to. The honest framing is that the increments $2,8,8,18,18,32,32$ each
occur twice except the first, and repeats need no explanation: the question is at the
\emph{growth points} --- why an increment grows after exactly two uses, and why to the next value
rather than another repeat. After two octets argon stands at $18$; a third octet would give $26$,
and $36$ is observed. After two $18$s xenon stands at $54$; a third would give $72$, and $86$ is
observed --- and $72$ is itself a capacity, so a further repeat is not absurd, it merely does not
happen. The orbital account obtains the doubling from the Madelung $(n{+}\ell)$ rule, which is
empirical, and the $(n{+}\ell)$ groups close at the alkaline earths rather than the noble gases;
so the problem is open in both frameworks, and we claim no more here than that it is open in ours.

\section{What closure requires}\label{sec:closure}

Closure is \textbf{exhaustion of a row's capacity}. At $Z = 2, 10, 18, 36, 54, 86, 118$ --- the
partial sums of $2, 8, 8, 18, 18, 32, 32$ --- every shell the period provides is full, so an
arriving charge finds no vacant site at that radius and must open a shell further out, which is
what begins the next period. Rubidium opening a shell outside a complete fourth row is the
mechanism, and it requires nothing about the shape of the outer layer.

This distinction matters, because the more familiar statement is false beyond the third row.
In rows 2 and 3 closure can equally be stated as a completed outermost octet, since there the row
\emph{is} one octet and the two statements coincide; and in those rows no atom can carry a complete
outermost octet without being a noble gas. From the fourth row on the two part company: the row is
a pair and two octets, and the pair may stand outside. Krypton's outermost eight is not $4{+}4$
under either filling rule we use, and in the orbital account $[\mathrm{Ar}]3d^{10}4s^24p^6$ is a
pair and six, not an octet. Scandium through zinc all carry a completed outermost \emph{pair}
without being noble gases --- ten counterexamples in a row. So it is capacity exhaustion, not the
outer layer's shape, that carries the account.

\section{The octet as two radial shells}\label{sec:subshells}

A filled octet is realised over \emph{two} radii rather than as a single eight-shell. The
diagnostic is energetic: forcing the eight onto one shell puts the charge closer to the kernel and
over-binds by $9$--$22\%$ across the atoms where it can be tested, which is a large and systematic
effect, not a marginal preference. Neon is representative: at $2{+}4{+}4$ the computed total is
$-131.731$~Ha against a reference $-128.938$~Ha, while a single outer eight collapses far below
observation.

Three qualifications belong with that result, and all three were added or sharpened in review.

\paragraph{The over-binding shows two shells, not which two.}
That a single eight-shell over-binds establishes only that the octet occupies two radial shells.
It does \emph{not} discriminate $4{+}4$ from $2{+}6$, which is also two shells. What selects four
is that cap-4 is the best performer among the candidate capacities $2, 4, 6, 8$ when the whole
table is computed under each.

\paragraph{The decomposition is not new.}
Reading the octet as two interpenetrating tetrahedral sets of four is \textbf{Linnett's double
quartet}~\cite{Linnett1961}, which recasts Lewis's cubical octet in exactly that way, and the idea
remains current. It should not be presented as novel here. What is different is the route:
Linnett separates the two tetrahedra by \emph{spin}, with Fermi correlation within a quartet and
Coulomb repulsion between quartets. RealQM has no spin at all and arrives at the same arrangement
from Coulomb interaction and non-overlap alone. A spin-free derivation of Linnett's decomposition
is the defensible contribution, and it is a better claim than priority over the geometry.

\paragraph{The cube is not a Coulomb minimum, and is not solver output.}
Two cautions. First, for eight point charges on a sphere the unique Coulomb minimiser is the
\emph{square antiprism}, not the cube, by $0.33\%$ --- $19.6753$ against $19.7408$ --- and this is
a theorem~\cite{Thomson8}. So the cube must never be offered as a Coulomb minimum. Two things
limit the damage: the antiprism is itself a $4{+}4$, two squares twisted by $45^\circ$, so the
theorem refutes the two-\emph{tetrahedra} reading of a single sphere rather than the $4{+}4$
decomposition; and the model never places eight charges on one sphere, the octet being realised
over two radii, which a single-sphere theorem does not reach. Second, the spherically-homogenised
solver used for the table below carries \textbf{no angular coordinate}: its shells are disjoint
radial intervals, so a configuration written $2{+}4{+}4{+}4{+}4$ is five radial shells and the
cube cannot be represented in it at all. Argon's measured shell boundaries are uniformly nested,
with successive mid-radius ratios $2.65$, $2.29$, $2.16$, and not paired as a cube reading would
predict. The cube is therefore an argument about counts and symmetry, not a computed result. The
capacities, the parity rule and the closures are arithmetic about counts and are unaffected.

\section{Validation: hydrogen to krypton}\label{sec:validation}

Table~\ref{tab:atoms} gives the computed total energies through argon, and
Table~\ref{tab:atoms4} the fourth row. Both are all-electron on the bare kernel with nothing
adjusted per element. The mesh is chosen by $Z$ rather than held fixed: the innermost shell
contracts as $r_{1s} \approx 4.3/Z$, so a single mesh cannot serve the whole table, and
$N = 400$ through argon with $N = 800$ from potassium keeps that shell spanned by a comparable
number of cells at both ends.

\begin{table}[h]
\centering
\small
\renewcommand{\arraystretch}{1.08}
\caption{Level-1 total energies, balanced cap-4 packing on the spherically-homogenised solver,
against near-exact nonrelativistic references~\cite{Chakravorty}. Uniform mesh $N=400$.
Mean $|\text{error}| = 1.6\%$, maximum $4.5\%$ (Si, with S at $4.3\%$); 13 of 17 within $2.5\%$.
The reference column is not measured data --- a total electronic energy is not an observable, and
these values are composites of experiment and calculation.}\label{tab:atoms}
\begin{tabular}{llcrrr}
\toprule
$Z$ & atom & config & RealQM (Ha) & reference (Ha) & error \\
\midrule
2  & He & $2$                 & $-2.845$   & $-2.904$   & $-2.0\%$ \\
3  & Li & $2{+}1$             & $-7.522$   & $-7.478$   & $+0.6\%$ \\
4  & Be & $2{+}1{+}1$         & $-14.580$  & $-14.667$  & $-0.6\%$ \\
5  & B  & $2{+}2{+}1$         & $-24.707$  & $-24.654$  & $+0.2\%$ \\
6  & C  & $2{+}2{+}2$         & $-37.446$  & $-37.845$  & $-1.1\%$ \\
7  & N  & $2{+}3{+}2$         & $-54.737$  & $-54.589$  & $+0.3\%$ \\
8  & O  & $2{+}3{+}3$         & $-75.382$  & $-75.067$  & $+0.4\%$ \\
9  & F  & $2{+}4{+}3$         & $-102.485$ & $-99.734$  & $+2.8\%$ \\
10 & Ne & $2{+}4{+}4$         & $-131.731$ & $-128.938$ & $+2.2\%$ \\
11 & Na & $2{+}4{+}4{+}1$     & $-165.545$ & $-162.255$ & $+2.0\%$ \\
12 & Mg & $2{+}4{+}4{+}1{+}1$ & $-201.704$ & $-200.053$ & $+0.8\%$ \\
13 & Al & $2{+}4{+}4{+}2{+}1$ & $-249.405$ & $-242.346$ & $+2.9\%$ \\
14 & Si & $2{+}4{+}4{+}2{+}2$ & $-302.465$ & $-289.359$ & $+4.5\%$ \\
15 & P  & $2{+}4{+}4{+}3{+}2$ & $-347.534$ & $-341.259$ & $+1.8\%$ \\
16 & S  & $2{+}4{+}4{+}3{+}3$ & $-415.030$ & $-398.110$ & $+4.3\%$ \\
17 & Cl & $2{+}4{+}4{+}4{+}3$ & $-466.105$ & $-460.148$ & $+1.3\%$ \\
18 & Ar & $2{+}4{+}4{+}4{+}4$ & $-527.994$ & $-527.540$ & $+0.1\%$ \\
\bottomrule
\end{tabular}
\end{table}

\begin{table}[h]
\centering
\small
\caption{Fourth row, $Z=19$--$36$, all-electron on the $N=800$ mesh. Mean
$|\text{error}| = 3.3\%$, maximum $6.7\%$ (Se); \textbf{all eighteen entries are under-bound},
which is the direction the spherical reduction predicts. Configurations from scandium onward are
\emph{supplied} from the ionisation order and are not produced by the outward build-up: what the
row demonstrates is that the solver handles the nesting and returns the energies once a
configuration is given.}\label{tab:atoms4}
\begin{tabular}{llrrr}
\toprule
$Z$ & atom & RealQM (Ha) & reference (Ha) & error \\
\midrule
19 & K  & $-584.495$  & $-599.230$  & $-2.5\%$ \\
20 & Ca & $-666.030$  & $-676.758$  & $-1.6\%$ \\
36 & Kr & $-2663.335$ & $-2752.050$ & $-3.2\%$ \\
\bottomrule
\end{tabular}
\end{table}

\paragraph{Why no sign is read off these tables.}
Argon's $+0.1\%$ is the closest agreement in Table~\ref{tab:atoms} and it is not robust. On the
same domain, refining to $N = 724$ and $N = 1024$ --- where the $1s$ spans $17$ and $25$ cells
rather than the $10$ it spans at $N = 400$ --- moves argon to $-511.7$ and $-511.3$~Ha, some $3\%$
\emph{under}-bound, with the two finest values agreeing to $0.08\%$. Both are step-converged:
quadrupling the relaxation at fixed mesh moves the energy by $0.005$~Ha. So the near-exact argon
entry is an artefact of the production mesh and nothing rests on it. In particular, an earlier
claim that these errors have the opposite sign to Hartree--Fock has been withdrawn, since the sign
of argon's error is itself a function of the grid. What survives refinement, and what the
spherical reduction predicts, is under-binding.

These are therefore the errors \emph{of this mesh} and not a converged accuracy. Single entries
move by a few per cent under refinement, and the movement is systematic rather than random, so it
does not average down over seventeen atoms. The mean should be read at the same level as the
entries --- a few per cent --- which is what a radial solve with no basis sets, no antisymmetry
and no two-electron integrals delivers. The largest error is at silicon, a light atom, so the
deviation is not a monotone accumulation with $Z$ and should not be described as one.

\paragraph{What the per-cent figures measure.}
The reference column is itself a composite of experiment and calculation, built from a
construction that already contains standard quantum mechanics. So these percentages measure
distance from that construction and cannot be read as a direct verdict of one theory against
nature. The observables this model can be held to without that circularity are of a different
kind: electronegativity ordering, shell radii, and geometry.

\section{Ionization energies and the configuration energy}\label{sec:chi}

The ionization-energy material of an earlier draft has been withdrawn in its entirety, including
the identification of a computed eigenvalue with an ionization energy by a Koopmans-type argument.
The reason is structural rather than numerical and is set out in Chapter~\ref{ch:limits}: a
total energy is a bulk sum and is robust, whereas an ionization energy depends on a single
loosely-bound valence charge and is fragile, and in the radial reduction the octet is split
radially rather than angularly, which places the outer tetrahedron too far out.

What replaces it is Allen's configuration energy, the mean valence-electron binding energy, which
the per-electron energies make immediately available. Across period 2 the computed
\emph{ordering} is exact, Li through Ne without inversion. The \emph{gradient} is too steep by a
factor of $2.23$: a span of $51.1$~eV against the observed $22.9$. An exact ordering with a
badly-scaled gradient is a fair summary of where this model stands on valence observables, and it
is reported as such rather than as agreement.

\section{What the results mean}\label{sec:meaning}

\paragraph{What they are.}
Neutral atoms are computed all-electron from hydrogen to krypton, to a few per cent, with $Z$ as
the only supplied number and nothing fitted. Two filled packings appear and no third, and the row
capacities of the periodic table are sums of those two alone, with the parity rule and the closed
form $\mathrm{cap}(n)$ following from them and the closures $2,10,18,36,54,86,118$ coming out
right. The Pauli principle does not appear as a separate postulate: the non-overlap constraint
between electron territories does that work.

\paragraph{What they are not.}
The angular structure within a shell is only partially represented, the $s/p$ distinction being
replaced by a radial split. Magnetic moments, fine structure and hyperfine structure are absent.
Ionization energies are outside range. Two further limits are worth stating plainly because they
bound what the construction can claim:
\begin{itemize}
\item \textbf{Capacities are derived at every $Z$; the filling is not, beyond carbon.} The
outward build-up derives lithium, beryllium, boron and carbon. Nitrogen through calcium use the
balanced cap-4 rule, adopted because it reproduces the known configurations; from scandium the
ionisation order is imported. So the energies are computed in configurations the model did not
itself choose, and closing that gap is what would make this an account of the table rather than a
calculation performed on one.
\item \textbf{Period doubling is unexplained}, here as in the orbital account.
\end{itemize}
An earlier claim that the build-up places $4s$ before $3d$ has been withdrawn: the $s^2$ shell was
placed outermost by construction, so the run tested its own input.

\section{Implications for the rest of Part III}\label{sec:to-part-iii}

The Level-1 agreement at a few per cent across hydrogen to krypton fixes the bottom of the
hierarchy. From here the higher-level architectures are derived by reduction from Level 1 and
checked against observation in turn.

\begin{itemize}
\item Chapter~\ref{ch:h2} works out the H$_2$ covalent bond at Level 1 with no kernel softening,
and recovers the experimental dissociation energy to $3\%$.
\item Chapter~\ref{ch:hydrides} sweeps Level-3 architectures across periodic-table groups, with
kernel parameters $(Z_{\rm kernel}, r_c)$ calibrated against the Level-1 valence sector.
\item Chapter~\ref{ch:dimers} tests the resulting Level-3 model against the S66 H-bond benchmark
and against reactive chemistry.
\item Chapter~\ref{ch:folding} stresses Level 4 with universal H-bond bias on polypeptide chains.
\item Chapter~\ref{ch:condensed} stresses Level 3 against three condensed-phase systems and
quantifies what is missing.
\end{itemize}

The thread connecting these chapters is the Level-1 anchor of this chapter. One caution carries
forward with it: the anchor is good to a few per cent and not better, the errors are systematic
rather than random so they do not average down, and a discrepancy at a higher level that is
smaller than a few per cent of the atomic total cannot be attributed to the higher level at all.

% --- End of Chapter 6 ---
